\section{Discussion and Conclusion}
We introduced \method, a lightweight, training-free controller for multi-agent LLM systems that addresses three persistent challenges in multi-agent collaboration: dynamic task routing, long-horizon credit assignment, and cold-start inefficiency. By combining Thompson-sampling delegation, reflection-driven updates, and memory-aware priors, \method{} maintains saturated task success while improving routing efficiency and robustness with minimal overhead.

\paragraph{Key insights.} Our findings support a growing body of evidence that lightweight, interpretable mechanisms can rival or surpass resource-intensive pipelines. The probabilistic control formulation resonates with models of agent trust~\cite{Wu21} and bandit methods for online decision-making~\cite{Chapelle11}, while avoiding the instability and sample inefficiency of deep RL~\cite{Mnih15}. The recursive loop mirrors ideas from search-based reasoning such as Monte Carlo Tree Search~\cite{Browne12}, yet operates at a fraction of the computational cost. In ablations, removing belief-guided routing increased poor-agent usage by 17.2\% and wasted attempts, underscoring the critical importance of online feedback for efficient coordination.

\paragraph{Limitations.} \method{} depends critically on judge reliability; biases in machine adjudication are well documented~\cite{Karpinska21} and can distort competency updates, potentially entrenching suboptimal routing decisions. Sequential recursion can add latency in long tasks, limiting applicability in real-time scenarios. Early cold-start behavior resembles random delegation before sufficient evidence accumulates to differentiate agent competencies. These limitations highlight the need for robust judge calibration, potential parallelization of candidate generation, and more informative prior initialization strategies.

\paragraph{Broader implications.} The transparency of \method{} is advantageous for responsible AI deployment: Beta posteriors, judge verdicts, and natural language rationales create an auditable decision trail~\cite{Doshi17}, contrasting sharply with opaque RL policies or heavily fine-tuned black-box controllers. However, adaptive down-weighting based on early evidence could prematurely exclude competent agents, particularly in scenarios with temporary performance fluctuations. Periodic recalibration and ensemble judge mechanisms are therefore important for reliable deployment in sensitive domains.

\paragraph{Future directions.} Several extensions warrant investigation. Hybridizing Bayesian delegation with model-based reinforcement learning could combine fast online adaptation with long-term strategic planning~\cite{Silver17}. Expanding beyond binary success/failure feedback into richer error taxonomies (e.g., reasoning errors, factual inaccuracies, incomplete coverage) would enable more fine-grained credit assignment and targeted agent improvement. Parallel candidate generation could significantly reduce latency while maintaining bounded computational budgets. Finally, programmatic verifiers and retrieval-grounded evaluation models~\cite{min2023factscore} offer promising directions for improving judge reliability and reducing calibration overhead.

\paragraph{Conclusion.} Robust collaboration in LLM collectives does not require complex black-box architectures or extensive training. Simple probabilistic mechanisms---belief-guided delegation, calibrated reflection,
and memory-aware priors---can transform collections of independent agents into cohesive,
adaptive systems. This ``fast and frugal'' approach offers a scalable, interpretable, and practically deployable path forward for multi-agent LLM research and real-world applications.
